%%%%%%%%%%%%%%%%%%%%%%%%%%%%%%%%%%%%%%%%%
% Latin Explorer — Masters Final Paper
% XeLaTeX, Libertinus Serif, APA biblatex
%%%%%%%%%%%%%%%%%%%%%%%%%%%%%%%%%%%%%%%%%

\documentclass[11pt, a4paper]{article}
\raggedbottom

% ---- XeLaTeX font stack ----
\usepackage{fontspec}
\usepackage{unicode-math}
\setmainfont{Libertinus Serif}
\setsansfont{Fira Sans}
\setmathfont{Libertinus Math}

% ---- Layout ----
\usepackage[top=0.9in, bottom=1in, left=1in, right=1in]{geometry}
\usepackage{setspace}
\doublespacing
\setlength{\parindent}{0pt}
\setlength{\parskip}{6pt}

% ---- Colors ----
\usepackage{xcolor}
\definecolor{PrimaryColor}{HTML}{191970}
\definecolor{SecondaryColor}{HTML}{2980B9}

% ---- Section headings ----
\usepackage[raggedright]{titlesec}
\titleformat{\section}
  {\rmfamily\Large\bfseries\color{PrimaryColor}}{\thesection}{1em}{}
\titleformat{\subsection}
  {\rmfamily\large\bfseries}{\thesubsection}{1em}{}
\titleformat{\subsubsection}
  {\rmfamily\normalsize\bfseries}{\thesubsubsection}{1em}{}
\titlespacing*{\section}{0pt}{0.4\baselineskip}{0.2\baselineskip}
\titlespacing*{\subsection}{0pt}{0.4\baselineskip}{0.2\baselineskip}

% ---- Bibliography ----
\usepackage{csquotes}
\usepackage[style=apa, backend=biber]{biblatex}
\addbibresource{references.bib}

% ---- Hyperlinks ----
\PassOptionsToPackage{hyphens}{url}
\usepackage[hidelinks]{hyperref}
\urlstyle{same}

% ---- Miscellaneous ----
\usepackage{microtype}
\usepackage{enumitem}
\setlist[itemize]{topsep=0pt, partopsep=0pt, parsep=0pt, itemsep=0pt}

\hyphenation{Eth-no-math-e-mat-ics eth-no-math-e-mat-ics}

% ---- Author / title metadata ----
\newcommand{\studentname}{[Student Name]}
\newcommand{\shorttitle}{Latin Explorer}
\newcommand{\institution}{Miami University, Oxford, Ohio}

% ---- Running headers ----
\usepackage{fancyhdr}
\pagestyle{fancy}
\fancyhf{}
\fancyhead[L]{\small\studentname}
\fancyhead[R]{\small\shorttitle}
\fancyfoot[C]{\thepage}
\renewcommand{\headrulewidth}{0.4pt}
\renewcommand{\footrulewidth}{0pt}
\setlength{\headheight}{14pt}

% ---- Custom title block ----
\makeatletter
\renewcommand{\@maketitle}{%
  \newpage\null
  \begin{center}
    {\LARGE\bfseries\color{PrimaryColor} \@title \par}
    \vspace{0.8em}
    {\large \studentname \par}
    \vspace{0.2em}
    {\normalsize \institution \par}
  \end{center}
  \par\vspace{0.5em}
}
\makeatother

\title{\textit{Latin Explorer}: Technology, Artificial Intelligence,\\
       and the Living Latin Classroom}
\author{}
\date{}

%% ============================================================
\begin{document}
%% ============================================================

\maketitle
\thispagestyle{plain}

\begin{center}
\itshape ``There is no royal road to Geometry.'' \upshape --- Euclid
\end{center}

\vspace{0.5em}

\begin{center}
\small
\href{https://lamkalm.github.io/latin-explorer/}{lamkalm.github.io/latin-explorer}
\quad\textbullet\quad
\href{https://github.com/lamkalm/latin-explorer}{github.com/lamkalm/latin-explorer}
\end{center}

\vspace{1em}

%% ============================================================
\section{Introduction}
%% ============================================================

On paper, the goal was to create a website to host interactive tools for
learning Latin as a living language. This website, presentation, and essay
represent the culmination of two years spent in a professional training
program. Going back to college in order to become a high school teacher did
not always make sense. One might as well learn how to swim by fishing---or
learn a language through a textbook. Enter Latin.

One of the perks of majoring in Latin Education as opposed to English or
Math is seeing how people react to learning that people go to school for
this. Their responses range from polite interest and delighted surprise to
mild amusement, thinly veiled condescension, or outright laughter. These
reactions are understandable to a point: Latin is quite possibly the driest,
most academic subject imaginable. The fact that many Latin programs---along
with many other non-technical, non-ELA and Math subjects---are dying seems
to support some of these reactions. On the other hand, studying the classics
often goes hand in hand with elitist attitudes (such attitudes extended to
eugenics in the early 20th century in places like Princeton)
\autocite[see, e.g.,][]{Stray1998}.

Language acquisition theories emphasize the importance of interaction,
context, and meaningful communication \autocite{Krashen1982}. As
\textcite{Yuldasheva2026} notes, ``while Latin is not commonly spoken,
modern approaches attempt to simulate communicative environments through
digital tools and interactive platforms.'' As a Latin student, I found
myself thinking multiple times: \textit{AI was built for this!} If I were
to settle on conversing with and learning from AI, how could I ever fairly
blame my students for getting up from their desks and going back home to do
the same?

Teaching Latin is unlike teaching Spanish, French, or any other modern
language. Outside of the Vatican and very few classrooms, Latin is
practically dead \autocite{Beard2013}. The focus shifts from writing and
speaking to decoding ancient texts. The sort of topical vocabulary taught in
other foreign language classrooms---weather, colors, emotions---are not as
useful as studying core vocabulary like prepositions, adverbs, and a litany
of verbs and nouns related to ancient Rome (e.g.\ \textit{servus, bellum,
ferrum}). How can teachers and students reasonably expect to study two sets
of material for one subject at the same time? Bridging these two approaches
was the primary aim of this project.

While most will concede that Latin was the language of Western scholarship
and the learned professions of medicine and law \autocite{Waquet2001}, a
piece from \textit{The Atlantic} makes a similar point: ``At Harvard
University, for example, the wide extension of the elective system led to
the abandonment many years ago of the requirement of Latin in college for
the degree of Bachelor of Arts.'' The article was written in 1917.

If the teaching and studying of the classics is anything more than
pretentious elitism, anyone choosing to seriously pursue this tradition
should not be bothered. The ancients themselves were oftentimes equipped to
meet arrogant ignorance with their own.

Using technology to teach something as academic or obsolete as Latin might
seem like a contradiction. One of my host teachers noticed that I was using
ChatGPT to help me with my Latin homework. His attitude towards me chilled
somewhat, and the next day he made a pointed comment about how filthy
cheating is. (I was having it generate definitions of words I did not know
and asking it questions about my translation on an ungraded assignment---not
a violation of academic integrity.) I chose not to respond emotionally, let
alone verbally. What might this kind of reaction teach us about Latin and AI?

The assumption that a student using AI must be cheating appears to be rooted
in mistrust of this emerging technology \autocite{Lim2023}. A famous Greek
mathematician once told his pupil, a king, that there is no royal road to
geometry. If I witnessed one of my students using AI in the classroom, I
would be no less skeptical---but I would ask questions before I accused.

%% ============================================================
\section{Artificial Intelligence}
%% ============================================================

My scholastic journey in the post-reality era would not be what it is---or
is not---without AI. \textcite{Khan2023} expressed the following in a 2023
TED Talk:
\begin{quote}
We're at the cusp of using AI for probably the biggest positive
transformation that education has ever seen [\ldots] The way we're going to
do that is by giving every student on the planet an artificially intelligent
but amazing personal tutor, and we're going to give every teacher on the
planet an amazing, artificially intelligent teaching assistant.
\end{quote}
This aspiration echoes \textcite{Bloom1984}'s famous ``two-sigma problem,''
which demonstrated that one-on-one human tutoring produces learning gains
two standard deviations above conventional classroom instruction---gains
that no scalable intervention had yet achieved. AI-powered tutoring is now
the most promising candidate for closing that gap \autocite{Mollick2023}.

\textcite{Colace2025} writes in ``Teaching Latin, inclusion, and Artificial
Intelligence: What do Teachers Think?'':
\begin{quote}
We live in an era where the digital environment is no longer merely an
external dimension of reality but increasingly becomes an integral and
structuring component, significantly impacting our personal, professional,
and social lives\ldots\ the use of adaptive learning systems, which open up
possibilities for constructing personalized pathways that respond to
learners' specific difficulties; the integration of gamification to
strengthen motivation through playful and immersive environments; the
reflective use of assisted translation to support linguistic and semantic
analysis of texts without replacing the student's hermeneutic activity; and
the adoption of specific digital tools that help reinforce reading and
comprehension practices, fostering interdisciplinarity and inclusivity.
\end{quote}

Computer-Assisted Language Learning (CALL) has a decades-long research
tradition documenting the effectiveness of digital tools in language
instruction \autocite{Warschauer1998, Chapelle2001}. \textit{Latin Explorer}
situates itself within this tradition while deliberately scoping its
ambitions to the tools most useful in a living-Latin context.

\textit{Latin Explorer} will focus exclusively on digital tools. I would
not have been able to build the site without the use of artificial
intelligence; however, my choice not to integrate AI into the site itself
was deliberate. Part of my journey as a Latin student included---and
continues to include---the use of AI. From generating definitions of
unfamiliar words to debating translation choices, AI has proved useful.
This is especially true when it comes to holding conversations in Latin,
since, as previously noted, opportunities to actually use the language are
vanishingly rare.

There is no room for AI in my Latin classroom---there must be limits and
boundaries. It would not have been possible for me to build such a site
without agentic coding, but this is where I draw my line. This position
aligns with calls in the literature for what \textcite{Colace2025} term the
``reflective use of assisted translation'': AI as scaffold, not substitute.

%% ============================================================
\section{Building the Applets}
%% ============================================================

I had very little coding experience before creating this project. I had
heard of things like GitHub and agentic coding before, but never knew what
they were. My advisor and teacher Todd Edwards was enormously helpful, and I
greatly appreciate his patience, understanding, and guidance. The two
programs I used were VS Code and GitHub. Just as students might feel
intimidated by solving a math problem, I was oftentimes dissuaded from
fixing errors in my code. When GitHub one day refused to update and VS Code
was no help at all, I researched some commands to run in my terminal, and it
finally worked. I am not a computer person, so I rediscovered that special
feeling of overcoming an insecurity. I hope that I will remember this when I
have a student struggling with Latin or getting nervous.

%% ------------------------------------------------------------
\subsection{Virgil's \textit{Georgics}}
%% ------------------------------------------------------------

The reason I chose to learn Latin was to be able to read Virgil in the
original. It seemed fitting to choose Virgil as the author I would feature
in \textit{Latin Explorer}. Translating an ancient Roman farming almanac is
just one of the joys I will be able to share with my future students.

The \textit{Georgics} applet allows students to submit questions that the
teacher can review and incorporate into a lesson. Student responses are
submitted online; they are able to ask specific questions about grammar,
vocabulary, and the choices they face as translators. The teacher can not
only export each student's translation for grading, but see at a glance
where in the text students are struggling. Reading ancient texts with
attention to language and meaning is at the heart of what scholars call
``deep reading'' \autocite{Wolf2018}, and this applet is designed to
support rather than shortcut that process.

%% ------------------------------------------------------------
\subsection{A Living Language: Graffiti and Etymology}
%% ------------------------------------------------------------

The goal of this applet is to connect students' knowledge of English to
Latin. Latin's influence on English was primarily lexical: at least 50\% of
English vocabulary is derived from Latin \autocite{Ostler2007}. This is
oftentimes one of the most effective ways to engage students' interest in
the subject. With this applet, students are encouraged to search for
derivatives in whatever languages they speak, connecting the unit topic to
their own linguistic experience---an approach consistent with
heritage-language pedagogy \autocite{Polinsky2018}.

The graffiti applet encourages students to communicate in an interactive
way. The ability to write Latin is often not taught in the classroom, yet
writing reinforces the morphological and syntactic patterns students learn
through rote memorization in a much more meaningful way
\autocite{SharwoodSmith1993}. Topics such as colors, feelings, and weather allow
students to exercise vocabulary while tracking their own progress since the
start of the year. Prompts related to a unit topic are provided to facilitate
mastery of sentence frames and vocabulary.

%% ============================================================
\section{What I Learned at Miami}
%% ============================================================

\textit{Latin Explorer} was never intended to be used in a classroom. Dr.\
Edwards suggested that developing applets for helping students to learn
Latin would not only demonstrate my growth to my defense committee, but
would also be a way to share what I had learned about AI and basic coding
with other foreign language teachers. He also suggested that learning
agentic coding might serve me well outside of Miami and beyond teaching.
Using AI to help create this site and learn my content area has revealed its
seemingly infinite potential. At the very least, I am now better equipped to
understand my students' choice to use AI than I would have been otherwise.

There are several talking points used to pitch a subject most consider
obsolete, most of which are as flawed as they are dated. The most persuasive
line of reasoning came to me one day when I learned that a group of students
were having lunch in a teacher's classroom. ``They're learning Greek.''
``I thought you said Walnut Hills cut Greek?'' ``I did---they're doing it for
fun.''

The heart and mind are sprawling networks---a throne apiece. From the
humblest nerve to the very seat of our native royalty, pools, corridors,
courts, and halls where humanity---what it means to be alive---reside.
Learning is about helping our students make connections. It is my hope that
Latin, a subject very few will ever consider central to their lives, might
allow our students to unlock themselves, and come to a knowledge of the most
high---to make the ultimate connection between themselves and the most high.
What they choose to do from there can only ever be a choice they must make
for themselves. What we are teaching them is how to make as informed a
decision as they can. Latin is one path among many, and it is our job to
prepare students for their journey.

%% ============================================================
\printbibliography
%% ============================================================

\end{document}
